
\subsection{感应定律}

\begin{tabularx}{\columnwidth}{XX}
  磁通量&$\Phi=BS\cos\theta$\\
  法拉第电磁感应定律&$E=n\left|\dfrac{\Delta\Phi}{\Delta t}\right|$\\
  感应电流&$I=\dfrac{E}{R+r}$\\
  回路电荷量&$q=\dfrac{n\Delta\Phi}{R+r}$\\
  导体棒切割&$E=Blv$\\
\end{tabularx}

\begin{formulanotebox}
法拉第定律给大小，楞次定律判方向。只算电动势大小而不判断方向，后续列回路方程和受力方程时容易符号错误。

$\Phi=BS\cos\theta$ 中的 $\theta$ 是磁场方向与面积法线的夹角，不是磁场与线圈平面的夹角。
\end{formulanotebox}

\textbf{楞次定律}

感应电流的磁场总要阻碍引起感应电流的磁通量变化。常用口诀如下：

\begin{itemize}[leftmargin=1.5em]
  \item 来拒去留：磁体靠近线圈时排斥，远离线圈时吸引。
  \item 增缩减扩：磁通量增大时回路有缩小趋势，磁通量减小时回路有扩张趋势。
  \item 增反减同：原磁场增强时感应磁场与原磁场相反，原磁场减弱时感应磁场与原磁场相同。
\end{itemize}

\begin{keypointbox}[title={\strut\textcolor{white}{\bfseries 重点知识点：感应电流方向}}]
判断方向的固定步骤是：先判断穿过回路的原磁通量如何变化，再确定感应磁场方向，最后用右手定则判断感应电流方向。

楞次定律中的“阻碍”不是阻止变化，而是使变化变慢；磁通量仍会按照外界原因增大或减小。
\end{keypointbox}

\begin{casetypebox}[title={\strut\textcolor{white}{\bfseries 常见题型：动生与感生电动势}}]
\begin{itemize}[leftmargin=1.5em]
  \item 导体或回路在磁场中运动产生的电动势常称动生电动势，典型公式为 $E=Blv$。
  \item 回路不动而磁场变化产生的电动势常称感生电动势，优先用 $E=n\left|\frac{\Delta\Phi}{\Delta t}\right|$。
  \item 无论动生还是感生，本质都是磁通量变化，方向都由楞次定律判断。
  \item 题目若问通过回路的电荷量，优先使用 $q=\frac{n\Delta\Phi}{R+r}$。
\end{itemize}
\end{casetypebox}
